%!TEX root = main.tex

%In this section, we show how to compute the pre-images of the string operations in Equation~(\ref{eqn-str-enc}). The operations $\myfilter_e$,  $\mysplit_u$, $\mymatch_e$, $\mymatchall_e$ and $\myjoin_u$ can be easily encoded in transducers, the computations of their pre-images can be derived by the work \cite{atva2020}. (The details can be found in Appendix~\ref{???}.) We focus on $\mywrite$, $\subseq$ and $\elem$. 

\cite{atva2020} introduced cost-enriched finite automata (CEFA for short) and showed that the pre-images of all the considered string operations involving the string and integer data types only can be effectively computed. 
% 
%The decision procedure in~\cite{atva2020} introduced cost-enriched finite automata (CEFA for short) and show that the pre-images of all the string operations involving the string and integer data type can be effectively computed.
The OSTRICH %string constraint solving 
framework 
%If all the string operations in a 
essentially informs that once the pre-image is computable, the straight-line %string constraint satisfy this property, then the constraint 
fragment can be solved by propagating the constraints encoded in CEFA  %of the form $\svaru \in \aut$ 
backwards. Our decision procedure for  $\slexstrlogic$ constraints follow this approach. 

We first recall the definition of CEFAs.
%Recall the cost-enriched finite automata (CEFAs~\cite{atva2020}). 
Intuitively, CEFAs add write-only cost registers to finite state automata, where ``write-only'' means that the cost registers can only be written/updated but cannot be read, i.e., they cannot be used in the guards of the transitions. 
%Moreover, an LIA \tl{remember to introduce this earlier} formula on the values of the registers is introduced as an accepting condition. 
%
%Recall that we assume that $\Sigma$ is the alphabet of the strings/elements in the sequences and $\separator \not \in \Sigma$ is a separator used to encode string sequences as strings. 


%\begin{definition}[Cost-Enriched Finite Automaton]
A CEFA $\aut$ over $\Sigma'$ is a  tuple $(R, Q, \Sigma', \delta, I, F)$ where
	\begin{itemize}
		\item $R = \{r_1, \cdots, r_k\}$ is a finite set of registers,
		\item $Q, I, F$ are the same as in {\nfa}, i.e., the set of states, the set of initial states, and the set of final states, respectively, 
		%    \item $R = (r_1\cdots r_n)$ is a vector of mutually distinct cost registers,
		\item $\delta \subseteq Q \times \Sigma' \times Q \times \Int^R$ is a transition relation, where $\Int^R$ denotes the updates on the values of registers. (Recall that $\Sigma':=\Sigma \cup \{\separator\}$ is the extended alphabet.)
		%  set whose elements are tuples $(q, c, q', \myvec{v})$ where $q, q'$ are states of $Q$, $c$ is a letter in alphabet $\Sigma\cup\{\epsilon\}$ and $\myvec{v}$ is the cost update function for registers, which is an integer vector whose $i$-th element is the incremental value of register $r_i$.  $(q, a, q', \myvec{v})$ is written as $q\xrightarrow[\myvec{v}]{a} q'$ for readability.
%		\item $acc \in \Phi(R)$ is an LIA formula specifying an accepting condition.
		%    a linear integer arithmetic constraint function on final states. $\theta$ is called \emph{accepting condition}.
	\end{itemize}
%\end{definition}
We write $R_\aut$ for the set of registers of $\aut$. %We assume a linear order on $R$ and write 
Typically, $R_\aut$ is represented as a vector $(r_1, \cdots, r_k)$. Accordingly, we write assignments $r_i:= r_i + v_i$ for all $i\in [k]$ %element of $\Int^R$ 
simply as 
%its values are represented as 
$\vec{v}=(v_1, \cdots, v_k)$, i.e., $r_i$ is to be incremented by $v_i$ for each $i \in [k]$. We also write a transition $(q, a, q', \vec{v}) \in \delta$ as $q \xrightarrow[\vec{v}]{a} q'$. In particular, $(q,a,q')$ stands for the transition $(q,a,q',())$ without any updates.
%The definition of CEFA above is slightly different from the definition in~\cite{atva2020} in the sense that accepting conditions on register values
%are attached to final states. 

%The semantics of CEFA is defined as follows. 
Let $\aut = (R, Q, \Sigma', \delta, I, F)$ be a CEFA. 
%A \emph{configuration} of $\aut$ is a pair $(q, \vec{v})$ where $q \in Q$ and $\vec{v}$ is a vector denoting the values of registers.  
%An \emph{initial configuration} of $\aut$ is $(q_0, \vec{0})$ with $q_0 \in I$, where the value of each register is zero. 
A \emph{run} of $\aut$ on a string $w = a_1 \cdots a_n$ over $\Sigma'$ is a sequence $q_0 \xrightarrow[\myvec{v_1}]{a_1} q_1 \cdots q_{n-1}\xrightarrow[\myvec{v_n}]{a_n} q_n$ such that $q_0 \in I$ and $q_{i-1} \xrightarrow[\myvec{v_i}]{a_i} q_i$ for each $i \in [n]$. 
%
%The run $q_0 \xrightarrow[\myvec{v_1}]{a_1} q_1 \cdots q_{n-1}\xrightarrow[\myvec{v_n}]{a_n} q_n$ 
It is \emph{accepting} if $q_n \in F$, and 
%
the vector $\myvec{c}= \sum_{j \in [n]} \myvec{v_j}$ is defined as the \emph{cost} of an accepting run. %$q_0 \xrightarrow[\myvec{v_1}]{a_1} q_1 \cdots q_{n-1}\xrightarrow[\myvec{v_n}]{a_n} q_n$. 
(Note that all registers are initialized to zero.) 
%
%and updated to $\sum \limits_{j \in [n]} \myvec{v_j}$ after all the transitions in the run are executed. 
We write $\myvec{c} \in \aut(w)$ if there is an accepting run of $\aut$ on $w$ whose cost is $\myvec{c}$.  

The semantics of a CEFA $\aut$, denoted by $\Lang(\aut)$, is defined as the set of pairs $\{(w, \myvec{c}) \mid  \myvec{c} \in \aut(w)\}$.
%
In particular, if $I \cap F \neq \emptyset$, then $(\varepsilon, \myvec{0}) \in \Lang(\aut)$. Moreover, the \emph{output} of a CEFA $\aut$, denoted by $\cefaout(\aut)$, is defined as $\{\myvec{c} \mid  \exists w.\ \myvec{c} \in \aut(w)\}$.

%%%%%%%%%%%%%%%%%%%%%%%%%%%%%%%%%%%%%%%%%%%%%%%%%%%%%%%%%%%%%%%%%%%%%%%%%%%%%%%%%
The following example illustrates the backward propagation. 

\vspace*{-2mm}
%\begin{example}
%Consider the equality $\svarv = f(\svaru_1, it_1, \svaru_2, it_2)$ (where $\svaru_1,\svaru_2$ are string variables and $it_1, it_2$ are integer terms) and the CEFA $\aut = (R, Q, \Sigma', \delta, I, F)$, where $R= \{r_1\}$. %(that is, $R$ contains only one register). 
%Suppose that $f^{-1}(\aut)$ is a finite collection of tuples $(\autb_1, \autb_2, t_1)$, where $\autb_1$ and $\autb_2$ are CEFAs such that $R_{\autb_1}=(r'_{1,1}, r^{(1)}_1)$ and $R_{\autb_2} = (r'_{2,1}, r^{(2)}_1)$ and $t_1$ is an LIA formula over $r^{(1)}_1$ and $r^{(2)}_1$. Intuitively, $r'_{1,1}$ and $r'_{2,1}$ are the fresh registers corresponding to the two integer parameters,  and $r^{(1)}_1, r^{(2)}_1$ are the fresh registers introduced in $\autb_1$ and $\autb_2$ respectively for restoring $r_1$. Then the backward propagation of the constraint $\svarv \in \aut$ 
%%with respect to the equality $\svarv = f(\svaru_1, it_1, \svaru_2, it_2)$ 
%wrt. $f$ 
%nondeterministically chooses a tuple $(\autb_1, \autb_2, t_1)$, removes $\svarv = f(\svaru_1, it_1, \svaru_2, it_2)$ from the string constraint, and adds $t_1 \wedge it_1 = r'_{1,1} \wedge it_2 = r'_{2,1}$ to it. \qed
%\end{example}
\begin{example}
	Consider the equality $\svarv = f(\svaru_1, it_1, \svaru_2, it_2)$ (where $\svaru_1,\svaru_2$ are string variables, $it_1, it_2$ are integer terms, and $f$ is a function) and the CEFA $\aut = (R, Q, \Sigma', \delta, I, F)$, where $R= \{r_1\}$. %(that is, $R$ contains only one register). 
	%Moreover, let $t = r^{(1)}_1 + r^{(2)}_1$ which intuitively describes how $r_1$ can be calculated from $r^{(1)}_1$ and $r^{(2)}_1$. 
	Then an $R$-cost-enriched pre-image of $\aut$ under $f$, denoted by $f^{-1}_{R}(\aut)$, is a pair $(\cerl, t)$ where $\cerl \subseteq (\Sigma')^\ast \times (\intnum \times \intnum) \times (\Sigma')^\ast \times (\intnum \times \intnum)$, $t$ is an LIA term over $\{r^{(1)}_1, r^{(2)}_1\}$, and $\cerl$ satisfies that for every $(v, n') \in \Lang(\aut)$, there is $(u_1, n_1, n'_1, u_2, n_2, n'_2) \in \cerl$ such that $f(u_1, n_1, u_2, n_2) = v$, and  $n'= t[n'_1/r^{(1)}_1, n'_2/r^{(2)}_1]$. 
	%
	We say that a pre-image $f^{-1}_{R}(\aut) = (\cerl, t)$ is CEFA-definable if $\cerl$ can be expressed as a finite collection of CEFA tuples $(\autb_{j, 1}, \autb_{j, 2})$ with $j \in [k]$ such that $R_{\autb_{j,1}}=(r'_{1,1}, r^{(1)}_1)$, $R_{\autb_{j, 2}} = (r'_{2,1}, r^{(2)}_1)$, and $\cerl = \bigcup_{j \in [k]} \Lang(\autb_{j, 1}) \times \Lang(\autb_{j, 2})$.  
	%Suppose $t = r^{(1)}_1 + r^{(2)}_1$ and $f^{-1}_{R, t}(\aut)$ is computed as a finite collection of tuples $(\autb_1, \autb_2)$, where $\autb_1$ and $\autb_2$ are CEFAs such that $R_{\autb_1}=(r'_{1,1}, r^{(1)}_1)$ and $R_{\autb_2} = (r'_{2,1}, r^{(2)}_1)$ and $t_1$ is an LIA formula over $r^{(1)}_1$ and $r^{(2)}_1$. Intuitively, 
	Intuitively, $r'_{1,1}$ and $r'_{2,1}$ are the fresh registers corresponding to the two integer parameters, and $r^{(1)}_1, r^{(2)}_1$ are the fresh registers introduced in $\autb_{j, 1}$ and $\autb_{j, 2}$ respectively for computing $r_1$ according to $t$. Then the backward propagation of the constraint $\svarv \in \aut$ 
	%with respect to the equality $\svarv = f(\svaru_1, it_1, \svaru_2, it_2)$ 
	under $f$ nondeterministically chooses a tuple $(\autb_{j, 1}, \autb_{j, 2})$, removes $\svarv = f(\svaru_1, it_1, \svaru_2, it_2)$ from the string constraint, and adds $\svaru_1 \in \autb_{j, 1} \wedge \svaru_2 \in \autb_{j, 2} \wedge r_1 = t \wedge it_1 = r'_{1,1} \wedge it_2 = r'_{2,1}$ to the string constraint.
	%removes $\svarv = f(\svaru_1, it_1, \svaru_2, it_2)$ from the string constraint, and adds $r_1 = t \wedge it_1 = r'_{1,1} \wedge it_2 = r'_{2,1}$ to the string constraint. \qed
\end{example}
\vspace*{-2mm}



%%%%%%%%%%%%%%%%%%%%%%%%%%%%%%%%%%%%%%%%%%%%%%%%%%%%%%%%%%%%%%%%%%%%%%%%%%%%%%%%%%%%%%%%%%%%%%%%%%%%%%% 

We now show how to compute the pre-images of $\mywrite$, $\subseq$, $\elem$ and $\myseqlen$. Note that the pre-image computation utilizes an NFA $\aut_0$ (resp. $\aut_1$) that formats the strings that represent the string sequences (resp. strings), namely, $\aut_0$ (resp. $\aut_1$) recognizes the language $\separator (\Sigma \separator)^*$ (resp. $\Sigma^*$). 

\paragraph*{Pre-image of $\mywrite$.}
Let $\aut = (R, Q, \Sigma', \delta, I, F)$ be a CEFA. 
Intuitively, $\mywrite^{-1}(\aut)$ is computed as the collection $((\autb_{(p, q)} \cap \aut_0, \aut_{R_2/R}[p, q] \cap \aut_1), t_{(p,q)})_{(p, q) \in Q \times Q}$, where  $\autb_{(p, q)}$, $ \aut_{R_2/R}[p, q]$ and $t_{(p,q)}$ are constructed as follows. 
\begin{itemize}
\item $\autb_{(p, q)}$ is constructed by the following procedure. 
\begin{itemize}
%\item We compute an existential LIA formula $\varphi_{p, q}(\myvec{x}, \myvec{y})$ to capture the reachability relation of $\aut[p, q]$. 
%
\item We create two copies of $R$, say $R_1, R_2$. Moreover, we introduce a fresh cost register, say $r'$, to count the number of occurrences of $\separator$. 
%
\item We run $\aut$ on the input string and increase the cost register $r'$ each time when $\separator$ is read. Moreover, the registers in $R_1$, instead of $R$, are updated in the transitions. 
%
\item When the current symbol is $\separator$, we nondeterministically choose to stop increasing $r'$ and pause the running of $\aut$ at the state $p$. Let us call this position as the pause position.
%
\item Finally, when reading the first $\separator$ after the pause position, we resume the run of $\aut$ at $q$, where the registers in $R_1$ (but not $r'$) are updated. 
%
\end{itemize}

\item $\aut_{R_2/R}[p, q]$ is obtained from $\aut[p, q]$ by replacing each register in $R$ with the corresponding copy in $R_2$, where $\aut[p, q]$ is the sub-automaton of $\aut$ that accepts the strings starting from the state $p$ and ending at the state $q$.
%
\item $t_{(p, q)} := (t_{(p, q), r})_{r \in R}$ and 
$t_{(p, q), r} = r^{(1)} + r^{(2)}$ for each $r \in R$.
%\item $t_{(p,q)} := \bigwedge_{r\in R} r = r^{(1)} + r^{(2)}$, where $r^{(1)}$ and $r^{(2)}$ are the copies of $r$ in $R_1$ and $R_2$ respectively.
\end{itemize}
Note that in the backward propagation of $\aut$ with respect to an equality $\svarz = \mywrite(\svarx, it, \svary)$, the constraint $r' = it $ is added to the LIA constraint to assert that $r'$ is equal to the index $it$. 
%
Formally, $\mywrite^{-1}(\aut)$ is a finite collection of $(\autb_{(p,q)}  \cap \aut_0, \aut_{R_2/R}[p, q] \cap \aut_1, t_{(p,q)})$, where $(p, q) \in Q \times Q$, and $\autb_{(p,q)} = (R_1 \cup \{r'\}, Q', \Sigma', \delta', I', F')$ such that  $Q' = Q \times \{pre, idle, post\}$, $I'  = I \times \{pre\}$, $F'  = F \times \{post\}$, and $\delta'$ comprises, 
\begin{itemize}
\item $((q_1, pre), a, (q_2, pre), (\myvec{v}, 0))$ such that $(q_1, a, q_2, \myvec{v}) \in \delta$ and $a \in \Sigma$ (thus $a \neq \separator$), 
%
\item $((q_1, pre), \separator, (q_2, pre), (\myvec{v}, 1))$ such that $(q_1, \separator, q_2, \myvec{v}) \in \delta$, 
%
\item $((q_1, pre), \separator, (p, idle), (\myvec{v}, 1))$ such that $(q_1, \separator, p, \myvec{v}) \in \delta$,
%
%\item $((q_1, idle), a, (q_1, idle), (\myvec{0}, 0))$ such that $a \in \Sigma$, 
%
%\item $((q_1, idle), \separator, (q_2, post), (\myvec{v}, 0))$ such that $(q, \separator, q_2, \myvec{v}) \in \delta$,
%\item $((q_1, idle), \separator, (q_2, post), (\myvec{v}, 0))$ such that $(q, \separator, q_2, \myvec{v}) \in \delta$,
%
\item $((p, idle), a, (p, idle), (\myvec{0}, 0))$ such that $a \in \Sigma$, 
\item $((p, idle), \separator, (q_2, post), (\myvec{v}, 0))$ such that $(q, \separator, q_2, \myvec{v}) \in \delta$,
%
\item $((q_1, post), a, (q_2, post), (\myvec{v}, 0))$ such that $(q_1, a, q_2, \myvec{v}) \in \delta$ and $a \in \Sigma$,
%
\item $((q_1, post), \separator, (q_2, post), (\myvec{v}, 0))$ such that $(q_1, \separator, q_2, \myvec{v}) \in \delta$.
\end{itemize}
From the construction of $\autb_{(p,q)}$, we can observe that the pair $(p,q)$ should satisfy that in $\aut$, there is a $\separator$-transition into $p$ and a $\separator$-transition out of $q$. Moreover, $q$ should be reachable from $p$ according to the  construction of $\aut_{R_2/R}[p, q] \cap \aut_1$. 
% \textcolor{blue}{Did not see how `` $q$ should be reachable from $p$'' is ensured.}
%Moreover, $acc' = acc$. 

\vspace*{-2mm}
\begin{example}
Let us consider the constraint $\svarv = \mywrite(\svaru, 1, \svary)\ \wedge\ \svarv \in \separator (a^+ \separator)^*\ \wedge\ \mystrlen(\svarv) \ge 4$. Let $\aut = (\{r_1\}, \{q_0, q_1, q_2\}, \Sigma', \delta, \{q_0\}, \{q_1\})$ be the CEFA, where the register $r_1$ is used to record the string length and $\delta$ comprises the transitions $q_0 \xrightarrow[(1)]{\separator} q_1 \xrightarrow[(1)]{a}q_2 \xrightarrow[(1)]{a} q_2 \xrightarrow[(1)]{\separator} q_1$. 
Let us consider the pairs $(p, q)$ in $\aut$ satisfying that there is a $\separator$-transition into $p$ and a $\separator$-transition out of $q$, moreover, $q$ should be reachable from $p$.
It is not hard to observe that there is exactly one such pair, that is, $(q_1, q_2)$. 
Therefore, $\mywrite^{-1}(\aut)$ comprises exactly one tuple $(\autb_{(q_1, q_2)}\cap \aut_0, \aut_{R_2/R}[q_1,q_2] \cap \aut_1, (r^{(1)}_1+r^{(2)}_1))$ such that
%and $(\autb_{(q_2, q_2)}\cap \aut_0, \aut_{R_2/R}[q_2,q_2] \cap \aut_1, (r^{(1)}_1+r^{(2)}_1))$ such that  
%\begin{itemize}
%\item 
$\autb_{(q_1,q_2)} = (\{r^{(1)}_1, r' \}, Q', \Sigma', \delta', I', F')$ where 
\begin{itemize}
\item $Q' = \{q_0, q_1, q_2\} \times \{pre, idle, post\}$, $I' = \{(q_0, pre)\}$, $F' = \{(q_1, post)\}$, and 
\item $\delta' $ comprises the following transitions 
\begin{itemize}
\item $((q_1, pre), a, (q_2, pre), (1,0))$, $((q_2, pre), a, (q_2, pre), (1,0))$, 
%$$(q_0, pre)\xrightarrow[(1,1)]{\separator}(q_1, pre)\xrightarrow[(1, 0)]{a}(q_2, pre) \xrightarrow[(1, 0)]{a}(q_2, pre) \xrightarrow[(1, 1)]{\separator}(q_1, pre),$$ 

\item $((q_0, pre), \separator, (q_1, pre), (1,1))$, $((q_2, pre), \separator, (q_1, pre), (1,1))$, 

\item $((q_0, pre), \separator, (q_1, idle), (1, 1))$ (since $(q_0, \separator, q_1, 1) \in \delta$),
%
\item $((q_1, idle), a', (q_1, idle), (0, 0))$ such that $a' \in \Sigma$, 
%
%\item $((q_1, idle), \separator, (q_2, post), (\myvec{v}, 0))$ such that $(q, \separator, q_2, \myvec{v}) \in \delta$,
\item $((q_1, idle), \separator, (q_1, post), (1, 0))$  (since $(q_2, \separator, q_1, 1) \in \delta$),
%
\item $((q_1, post), a, (q_2, post), (1, 0))$ and $((q_2, post), a, (q_2, post), (1, 0))$,
%
\item $((q_2, post), \separator, (q_1, post), (1, 0))$.
%
%$$
%\begin{array}{l}
%(q_0, pre)\xrightarrow[(1,1)]{\separator}(q_1, idle)\xrightarrow[(0, 0)]{\Sigma}(q_1, idle)\xrightarrow[(1, 0)]{\separator}(q_1, post)\\ 
%\xrightarrow[(1, 0)]{a} (q_2, post)  \xrightarrow[(1, 0)]{a} (q_2, post) \xrightarrow[(1, 0)]{\separator} (q_1, post),
%\end{array}
%$$ 
%where $(q_1, idle)\xrightarrow[(1, 0)]{\separator}(q_1, post) \in \delta'$ because $q_2 \xrightarrow[(1)]{\separator} q_1 \in \delta$.   
\end{itemize}
\end{itemize}
To illustrate how $\autb_{(q_1, q_2)}$ and $\aut_{R_2/R}[q_1,q_2]$ work, we can see that $(\separator aa \separator a \separator, 6)$ is accepted by $\aut$, moreover,  for any $b\in\Sigma$,
\begin{itemize}
\item $(\separator ab \separator a \separator, 4, 1)$ is accepted by $\autb_{(q_1, q_2)}$, witnessed by the following sequence of transitions, 
\[
\begin{array}{l}
(q_0, pre) \xrightarrow[(1, 1)]{\separator} (q_1, idle) \xrightarrow[(0, 0)]{a} (q_1, idle) \xrightarrow[(0, 0)]{b} (q_1, idle) \\
\xrightarrow[(1, 0)]{\separator} (q_1, post) \xrightarrow[(1, 0)]{a} (q_2, post) \xrightarrow[(1, 0)]{\separator} (q_1, post),
\end{array}
\]
%
\item and $(aa, 2)$ is accepted by $\aut_{R_2/R}[q_1,q_2]$, witnessed by the following sequence of transitions, 
$
q_1 \xrightarrow[(1)]{a} q_2  \xrightarrow[(1)]{a} q_2
$.  
\end{itemize}
Thus, $\svaru$ and $\svary$  in $\svarv = \mywrite(\svaru, 1, \svary)$
can be $\separator ab \separator a \separator$ and $aa$, respectively.\qed

%
%\item $\aut_{R_2/R}[q_1,q_2]$, 
%\item $\autb_{(q_1, q_1)}$ is similar to $\autb_{(q_1,q_2)}$ and is obtained from $\autb_{(q_1,q_2)}$ by replacing the transition $(q_1, idle)\xrightarrow[(0, 0)]{\separator}(q_2, post)$ in $\delta'$ with $(q_1, idle) \xrightarrow[(0, 0)]{\separator}(q_2, idle)$.
%
%$\autb_{(q_1,q_1)} = (\{r^{(1)}_1, r' \}, Q', \Sigma', \delta',I', F')$ where $Q' = \{q_0, q_1, q_2\} \times \{pre, idle, post\}$, $I' = \{(q_0, pre)\}$, $F' = \{(q_1, post)\}$, and $\delta' $ comprises the following transitions 
%$$(q_0, pre)\xrightarrow[(1,1)]{\separator}(q_1, pre)\xrightarrow[(1, 0)]{a}(q_2, pre) \xrightarrow[(1, 0)]{a}(q_2, pre) \xrightarrow[(1, 1)]{\separator}(q_1, pre),$$ 
%$$
%\begin{array}{l}
%(q_0, pre)\xrightarrow[(1,1)]{\separator}(q_1, idle)\xrightarrow[(0, 0)]{\Sigma}(q_1, idle)\xrightarrow[(0, 0)]{\separator}(q_1, post)\\ 
%\xrightarrow[(1, 0)]{a} (q_2, post) \xrightarrow[(1, 0)]{a} (q_2, post) \xrightarrow[(1, 0)]{\separator} (q_1, post),
%\end{array}
%$$ 
%One tuple of $\mywrite^{-1}(\aut)$ is $(\autb_{(q_1,q_2)}\cap \aut_0, \aut_{R_2/R}[q_1,q_2] \cap \aut_1, true)$ where $\autb_{(q_1,q_2)}$ consist of transitions $(q_0,pre)\xrightarrow[(1)]{\separator}(q_1,pre)\xrightarrow[(0)]{a}(q_2,pre)\xrightarrow[(1)]{\separator}(q_3,pre)$ and $(q_0,pre)\xrightarrow[(1)]{\separator}(q_1, idle)\xrightarrow[(0)]{\Sigma}(q_1,idle)\xrightarrow[(0)]{\separator}(q_2,post)$.
%$\aut_{R_2/R}[q_1,q_2]$ is $q_1\xrightarrow{a}q_2$. After back propagation, we remove the constraint $\svarv = \mywrite(\svaru, 1, \svary)$ and add the constraints $r' = 1\wedge \svaru \in \autb_{(q_1,q_2)}\cap \aut_0\wedge \svary \in \aut_{R_2/R}[q_1,q_2] \cap \aut_1$, where $r'$ is the fresh register introduced in $\autb_{(q_1,q_2)}$. Solving the new constraints, we obtain a solution $\svaru = \separator a \separator$ and $\svary = a$.
\end{example}
\vspace*{-2mm}

%$\myfilter_e$: seen as a special case of transducers

\paragraph*{Pre-image of $\subseq$.}
%$\subseq$: from $\aut$, introduce two fresh cost registers $r_1, r_2$, increment $r_1, r_2$ each time when a $\separator$ is read, nondeterministically choose to stop incrementing $r_1$, start running $\aut$ until nondeterministically choosing to stop incrementing $r_2$, moreover, check whether $\aut$ accepts when stopping incrementing $r_2$. 
%
Let $\aut = (R, Q, \Sigma', \delta, I, F)$ be a CEFA. 
Then $\subseq^{-1}(\aut)$ is computed as $(\autb\cap \aut_0, t)$ such that $t := true$ and $\autb$ introduces two fresh registers $r'_1, r'_2$ to count the two numbers of occurrences of $\separator$ that correspond to the starting position and the length of the subsequence respectively. Moreover, $\autb$ simulates the run of $\aut$ on the substring representing the subsequence returned by $\subseq$. 
%
Formally, $\subseq^{-1}(\aut)$ is computed as $(\autb \cap \aut_0, true)$, where $\autb = (R \cup \{r'_1, r'_2\}, Q \cup \{pre, post\}, \Sigma', \delta', \{pre\}, \{post\})$ such that $pre, post$ are two fresh states denoting the fact that $\autb$ is reading a symbol before and after the subsequence respectively and $\delta'$ comprises the following tuples, 
\begin{itemize}
\item $(pre, a, pre, (\myvec{0}, 0,0))$ such that $a \in \Sigma$ (thus $a \neq \separator$),  
%
\item $(pre, \separator, pre, (\myvec{0}, 1,0))$,
%
\item $(pre, \separator, q, (\myvec{v}, 1,0))$ such that $(p, \separator, q, \myvec{v}) \in \delta$ for some $p \in I$, 
%
\item $(q, a, q', (\myvec{v}, 0, 0))$ such that $(q, a, q', \myvec{v}) \in \delta$ and $a \in \Sigma$, 
%
\item $(q, \separator, q', (\myvec{v}, 0,1))$ such that $(q, \separator, q', \myvec{v}) \in \delta$,
%
\item $(q, \separator, post, (\myvec{v}, 0,1))$ such that $(q, \separator, q', \myvec{v}) \in \delta$ for some $q' \in F$,
%
\item $(post, a, post, (\myvec{0}, 0,0))$ and $(post, \separator, post, (\myvec{0}, 0,0))$ where $a \in \Sigma$.
\end{itemize}
Note that in the backward propagation of $\aut$ with respect to an equality $\svarz = \subseq(\svarx, it_1, it_2)$, the constraint $r'_1 = it_1 \wedge r'_2 = it_2$ is added to the LIA constraint to assert that the value of $r'_1$ and $r'_2$ are equal to the beginning index and the length of the subsequence respectively.

\vspace*{-2mm}
\begin{example}
Let $\svarv = \subseq(\svaru, 2, 1) \wedge \svarv \in \separator (a^+ \separator)^* \wedge \mystrlen(\svarv) \ge 4$ and $\aut = (\{r_1\}, \{q_0, q_1, q_2\}, \Sigma', \delta, \{q_0\}, \{q_1\})$ be the CEFA, where the register $r_1$ is used to record the string length and $\delta$ comprises the transitions $q_0 \xrightarrow[(1)]{\separator} q_1 \xrightarrow[(1)]{a}q_2 \xrightarrow[(1)]{a} q_2 \xrightarrow[(1)]{\separator} q_1$. 
Then $\subseq^{-1}(\aut)$ is $(\autb \cap \aut_0, true)$ such that 
$$\autb = (\{r^{(1)}_1, r'_1, r'_2\}, \{q_0, q_1, q_2, pre, post\}, \Sigma', \delta', \{pre\}, \{post\})$$
where 
$\delta' $ comprises the following transitions 
\begin{itemize}
\item $(pre, a', pre, (0, 0,0))$ such that $a' \in \Sigma$, 
%
\item $(pre, \separator, pre, (0, 1,0))$,
%
\item $(pre, \separator, q_1, (1, 1,0))$ (since $(q_0, \separator, q_1, 1) \in \delta$), 
%
\item $(q_1, a, q_2, (1, 0, 0)$, $(q_2, a, q_2, (1, 0, 0))$, 
%\item $(q, a, q', (\myvec{v}, 0, 0))$ such that $(q, a, q', \myvec{v}) \in \delta$ and $a \in \Sigma$, 
%
\item $(q_2, \separator, q_1, (1, 0, 1))$, 
%$(q, \separator, q', (\myvec{v}, 0,1))$ such that $(q, \separator, q', \myvec{v}) \in \delta$,
%
\item $(q_2, \separator, post, (1, 0, 1))$ (since $(q_2, \separator, q_1, 1) \in \delta$ and $q_1$ is an accepting state in $\aut$), 
%$(q, \separator, post, (\myvec{v}, 0,1))$ such that $(q, \separator, q', \myvec{v}) \in \delta$ for some $q' \in F$,
%
\item $(post, a', post, (0, 0,0))$ and $(post, \separator, post, (0, 0,0))$ where $a' \in \Sigma$.
\end{itemize}
To illustrate how $\autb$ works, we can see that $(\separator a  \separator a a \separator, 4, 2, 1)$ is accepted by $\autb$, witnessed by the following sequence of transitions, 

$
pre \xrightarrow[(0, 1, 0)]{\separator} pre \xrightarrow[(0, 0, 0)]{a} pre \xrightarrow[(1, 1, 0)]{\separator} q_1 \xrightarrow[(1, 0, 0)]{a} q_2 \xrightarrow[(1, 0, 0)]{a} q_2  \xrightarrow[(1, 0, 1)]{\separator} post.
$\qed
%where $\autb$ is 
%$pre\xrightarrow[(1, 0)]{\separator}pre\xrightarrow[(0, 0)]{\Sigma}pre\xrightarrow[(1, 0)]{\separator}q_1\xrightarrow[(0, 0)]{a}q_2\xrightarrow[(0, 1)]{\separator}post\xrightarrow[(0, 0)]{\Sigma'}post$. 
%After back propagation, we remove the constraint $\svarv = \subseq(\svaru, 1, 1)$ and add the constraints $r'_1 = 2\wedge r'_2 = 3\wedge \svaru \in \autb \cap \aut_0$, where $r'_1$ and $r'_2$ are the fresh registers introduced in $\autb$. Solving the new constraints, we obtain a solution $\svaru = \separator\separator a \separator$ which represents the sequence $(\epsilon, a)$.
\end{example}
\vspace*{-4mm}

%$\mysplit_u$: seen as a special case of transducers

%$\mymatch_e$ and $\mymatchall_e$: seen as special cases of transducers


\paragraph*{Pre-image of $\elem$.}
%$\elem$: seen as special cases of transducers
Let $\aut = (R, Q, \Sigma, \delta, I, F)$ be a CEFA. 
The computation of $\elem^{-1}(\aut)$ is similar to that of $\subseq^{-1}(\aut)$. Nevertheless, the construction is slightly different since the output of $\elem$ is a string, instead of a sequence. Intuitively, $\elem^{-1}(\aut)$ is computed as $(\autb \cap \aut_0,t)$ where $t := true$, a fresh register $r'$ is introduced to count the number of occurrences of $\separator$ that corresponds to the position where the element is extracted. Formally, $\elem^{-1}(\aut)$ is computed as $(\autb \cap \aut_0, true)$ such that $\autb = (R \cup \{r'\}, Q \cup \{pre, post\}, \Sigma', \delta', \{pre\}, \{post\})$, where $\delta'$ comprises the following tuples, 
\begin{itemize}
\item $(pre, a, pre, (\myvec{0}, 0))$ such that $a \in \Sigma$ (thus $a \neq \separator$),  
%
\item $(pre, \separator, pre, (\myvec{0}, 1))$,
%
\item $(pre, \separator, p, (\myvec{v}, 1))$ such that $p \in I$, 
%\item $(pre, \separator, p, (\myvec{0}, 1))$ such that $p \in I$, 
%
\item $(q, a, q', (\myvec{v}, 0))$ such that $(q, a, q', \myvec{v}) \in \delta$ and $a \in \Sigma$ (thus $a \neq \separator$), 
%
\item $(q, \separator, post, (\myvec{0}, 0))$ such that $q \in F$,
%
\item $(post, a, post, (\myvec{0}, 0))$ and $(post, \separator, post, (\myvec{0}, 0))$ where $a \in \Sigma$.
\end{itemize}
Note that in the backward propagation of $\aut$ with respect to an equality $\svarz = \elem(\svarx, it)$, the constraint $r' = it$ is added to the LIA constraint to assert that the value of $r'$ is equal to the index $it$.

%%%%%%%%%%%%%%%%%%%%%%%%%%%%%
\hide{
\begin{example}
Consider the constraints $\svarv = \elem(\svaru, 2) \wedge \svarv = a$. Then the automaton $\aut$ that $\svarv$ belongs to is $q_0\xrightarrow{a}q_1$. The $\elem^{-1}(\aut)$ is $(\autb \cap \aut_0, true)$, where $\autb$ is $pre\xrightarrow[(1)]{\separator}pre\xrightarrow[(0)]{\Sigma}pre\xrightarrow[(1)]{\separator}q_0\xrightarrow[(0)]{a}q_1\xrightarrow[(0)]{\separator}post\xrightarrow[(0)]{\Sigma'}post$. After back propagation, we remove the constraint $\svarv = \elem(\svaru, 1)$ and add the constraints $r' = 2\wedge \svaru \in \autb \cap \aut_0$, where $r'$ is the fresh register introduced in $\autb$. Solving the new constraints, we obtain a solution $\svaru = \separator\separator a \separator$ which represents the sequence $(\epsilon, a)$.
\end{example}
}
%$\myjoin_u$: special case of transducers

\medskip

Finally, we  remark that the operation $\myseqlen(\svarx)$ can be modeled as a CEFA, similarly to the CEFA for $\mystrlen$ in \cite{atva2020}.
%%%%%%%%%%%%%%%%%%%%%%%%%%%%%%%%%%
%%%%%%%%%%%%%%%%%%%%%%%%%%%%%%%%%%
\hide{
\paragraph*{Pre-image of $\myseqlen$.} 
Let $\aut = (R, Q, \Sigma, \delta, I, F)$ be a CEFA. 
The computation of $\myseqlen^{-1}(\aut)$ is similar to that of $\mystrlen^{-1}(\aut)$ in \cite{atva2020}, except that $\myseqlen^{-1}(\aut)$ count the number of occurrences of $\separator$ while $\mystrlen^{-1}(\aut)$ counts the length of the string. Formally, $\myseqlen^{-1}(\aut)$ is computed as $(\autb \cap \aut_0, t_{(p,q)})$ where $t_r = r^{(1)}$ for each $r \in R$, $\autb = (R \cup \{r'\}, \{q\}, \Sigma', \delta', \{q\}, \{q\})$ and $\delta'$ comprises the following tuples,
\begin{itemize}
\item $(q, a, q, (0))$ such that $a \in \Sigma$ (thus $a \neq \separator$),
\item $(q, \separator, q, (1))$.
\end{itemize}
Note that in the backward propagation of $\aut$ with respect to an equality $it = \myseqlen(\svaru)$, the constraint $r' = it$ is added to the LIA constraint to assert that the value of $r'$ is equal to the length of the sequence.

\begin{example}
	Let us consider the constraints $it = \myseqlen(\svaru) - 1\wedge it > 0$. The $\myseqlen^{-1}$ is $(\autb \cap \aut_0, true)$, where $\autb$ is $q\xrightarrow[(1)]{\separator}q\xrightarrow[(0)]{\Sigma}q$. After back propagation, we remove the constraint $it = \myseqlen(\svaru)$ and add the constraints $r' = it\wedge \svaru \in \autb \cap \aut_0$, where $r'$ is the fresh register introduced in $\autb$. Solving the new constraints, we obtain a solution $it = 1$ and $\svaru = \separator\separator$ which represents the sequence $(\epsilon)$. 
\end{example}
}



%% write by denghang ==================================================================
%% write by denghang ==================================================================
\hide{
	\subsection{Special CEFA for String Array}
	%%
	We define a special CEFA for string array. The CEFA uses a preserved character $``\arrseparator"$ as $array$ elements separator (note that $\arrseparator\not\in\Sigma$). The speciality of $\aut_s$ rely on  its construction.
	%%
	\paragraph{Array of fixed length and constant index} Considering an $array$ $A$ with fixed length $m$, we assume that each array element $A[i]$ $(i\in[1,m])$ is constrained by a CEFA $\aut_i\equiv (R_i, Q_i, \Sigma, \delta_i, I_i, F_i, \alpha_i)$. Note that the index of the array starts from 1.  Then we can construct the CEFA of the array $A$ as $\aut_s\equiv (R_s, Q_s, \Sigma\cup\{\arrseparator\}, \delta_s, I_s, F_s, \alpha_s)$, where:
	\begin{itemize}
		% \item $R_s = \bigcup\limits_{1\leq i \leq m} R_i$,
		\item $R_s = \bigcup\limits_{1\leq i \leq m} R_i\cup \myset{r}$ for a fresh register r,
		\item $Q_s = \bigcup\limits_{1\leq i \leq m} Q_i\cup \myset{q_0, q_f}$ for fresh states $q_0$ and $q_f$,
		\item $\delta_s$ is composed of four transitions set
		\begin{itemize}
			\item  $\{(q_0, \arrseparator, q_1, (\myvec{0}, 1))\mid q_1\in I_1\}$,
			\item  $\myset{(q_m, \sharp, q_f, (\myvec{0}, 0))\mid q_m\in F_m}$,
			\item  $\{(q, \arrseparator, q', (\myvec{0}, 1)) \mid q\in F_i, q'\in I_{i+1}, i\in[1,m]\}$ where $ I_{m+1}=\emptyset$,
			\item  and $\{(q, a, q', (\myvec{0},\myvec{v_i},\myvec{0}))\mid (q, a, q', \myvec{v_i})\in \aut_i, i\in[1,m]\}$,
		\end{itemize}
		where $(\myvec{0}, 1)$ is a vector that increases the value of $r$ by 1 and remains other registers unchanged, and $(\myvec{0},\myvec{v_i},\myvec{0})$ is a vector that updates the value of $R_i$ by $\myvec{v_i}$ and remains other registers unchanged. Similarly for $(\myvec{0}, 0)$ and so on.
		\item $I_s = \myset{q_0}$, $F_s = \myset{q_f},$ and $\alpha_s = \bigwedge\limits_{1\leq i \leq m} \alpha_i$.
		% and $\alpha_s = \bigwedge\limits_{1\leq i \leq m} \alpha_i$.
	\end{itemize}
	The resulting CEFA is shown in fig \ref{fig:fixed_len_cons_idx}, where $q_i\in I_i$ and $q_i'\in F_i$ for each $i\in [1,m]$. The main idea of the construction above is using $\arrseparator$ to separate each element of the array and one register to record the index of each element. We introduce two fresh states $q_0$ and $q_f$ to ensure each element is started and ended with at least one $\arrseparator $, which is beneficial to the construction of pre-images of string constraints (see \ref*{sec:pre_image} for more details). The accepting condition $\alpha_s$ is the conjunction of accepting conditions of each element.
	% It is obviously that the accepting word $w$ of $\aut_s$ is in language $\mathcal{L}(\aut_1)\arrseparator\mathcal{L}(\aut_2)\ ... \ \arrseparator \mathcal{L}(\aut_{m}) $.
	\begin{figure}[H]
		\centering
		\usetikzlibrary {automata, positioning, fit}
		\begin{tikzpicture}[initial text =, initial distance=3ex,
			node distance=1.2cm, auto,
			state/.style={circle, draw, inner sep=0pt, minimum size=5mm},
			accepting by double]
			
			\node[state,initial]            (q_0)                {$q_0$};
			\node[state]                    (q_1) [right=of q_0] {$q_1$};
			\node[state]                    (q_1') [right=of q_1] {$q_1'$};
			\node                           (omit)[right=of q_1'] {$\cdots$};
			\node[state]                    (q_m) [right=of omit] {$q_m$};
			\node[state]                    (q_m') [right=of q_m] {$q_m'$};
			\node[state, accepting]         (q_f)  [right=of q_m'] {$q_f$};
			
			\node[draw, fit=(q_1) (q_1'), inner sep=0.5em, label=above:$\aut_1$] [dashed] {};
			\node[draw, fit=(q_m) (q_m'), inner sep=0.5em, label=above:$\aut_m$] [dashed] {};
			
			\path[->]
			(q_0) edge [above] node {$\scriptstyle \arrseparator$} (q_1)
			(q_0) edge [below] node {${\scriptstyle r++}$} (q_1)
			(q_1') edge [above] node {$\scriptstyle \arrseparator$} (omit)
			(q_1') edge [below] node {${\scriptstyle r++}$} (omit)
			(omit) edge [above] node {$\scriptstyle \arrseparator$} (q_m)
			(omit) edge [below] node {${\scriptstyle r++}$} (q_m)
			(q_m') edge [above] node {$\scriptstyle \arrseparator$} (q_f)
			(q_m') edge [below] node {${\scriptstyle r}$} (q_f);
			
		\end{tikzpicture}
		\caption{The CEFA of $A$ with accepting condition $\alpha_s = \bigwedge\limits_{1\leq i \leq m} \alpha_i$}
		\label{fig:fixed_len_cons_idx}
	\end{figure}
	The construction of $\aut_s$ can be used to formalize the array containing constant string. For example, for a 2-length array $A$ constrained by $``a" = \myread(A, 1)\wedge ``b" = \myread(A, 2) \wedge \mylen(A) = 2$ , where $\myread(A, m)$ read the  element of $A$ at constant index $m$ and $\mylen(A)$ is the length of $A$, we can construct the automaton of $A$ as following: \\
	$(\emptyset, \{q_0,q_1,q_1',q_2,q_2'\}, \{a,b,\arrseparator\}, \{q_0\xrightarrow[(1)]{\arrseparator}q_1, q_1\xrightarrow[(0)]{a}q_1', q_1'\xrightarrow[(1)]{\arrseparator}q_2, q_2\xrightarrow[(0)]{b}q_2'\}, \{q_0\}, \{q_2'\}, \top)$. 
	%%
	\paragraph{Array of varible length and index} \todo{TODO}
}%% write by denghang ==================================================================
%% write by denghang ==================================================================



%%%%%%%%%%%%%%%%%%%%%%%%%%%%%%%%%%
%%%%%%%%%%%%%%%%%%%%%%%%%%%%%%%%%%
\hide{

\subsection{Pre-images of String Constraints}\label{sec:pre_image}

In this section, we define the pre-images of string constraint formulae. 
\pagebreak

We let $\aut_{n}$ be the automaton accepting all arrays with length $n$, that is, $\aut_{n} = q_0\xrightarrow{\arrseparator}q_1\cdots\xrightarrow{\arrseparator}q_{n}$ with $q_i\xrightarrow{\Sigma}q_i$ for $0\leq i \leq n-1$. $\aut_*$ is the automaton accepting all arrays with arbitrary length. Then the pre-images of operations are defined as follow.

\paragraph{- $i = \mylen(x)$:} As \cite{atva2020}.

\paragraph{- $x = y\cdot z$:} As \cite{ostrich}.

\paragraph{- $x = \transducer(y)$:} As \cite{ostrich}.

\paragraph{- $i = \mylen(A)$:} The pre-image for the array length function is the CEFA $\aut_{arrlen}=(R, Q, \Sigma\cup\{\arrseparator\}, \delta, I, F, \alpha)$, where
\begin{itemize}
	\item $R = \{r\}$ for a fresh register $r$, $Q = \{q_0, q_f\}$, $\delta = \myset{q_0\xrightarrow[(0)]{\Sigma}q_0, q_0\xrightarrow[(1)]{\sharp}q_0,  q_0\xrightarrow{\sharp}q_1} $, $I = \myset{q_0}$, $F = \myset{q_1}$, and $\alpha \equiv r = i$.
\end{itemize}

\paragraph{- $A = B[m:n]$:} Suppose that $A$ is constrained by $\aut_A=(R_A, Q_A, \Sigma\cup\{\arrseparator\}, \delta_A, I_A, F_A, \alpha_A)$, the pre-image of sub-array operation is defined as $\aut_m\cdot (\aut_A\cap \aut_{m-n})\cdot \aut_{*}$.

\paragraph{- $A = \transducer(B)$:} We consider $A$ and $B$ as string varibles and apply the pre-image computation in \cite{ostrich} directly.

\paragraph{- $x = \myread(A, m)$:} Suppose that $x$ is constrained by the CEFA $\aut_x=(R_x, Q_x, \Sigma, \delta_x, I_x, F_x, \alpha_x)$, the pre-image for the array read function is the CEFA $\aut_{read}=(R, Q, \Sigma\cup\{\arrseparator\}, \delta, I, F, \alpha)$ where
\begin{itemize}
	\item $R = R_x$,
	\item $Q = Q_x\cup\myset{q_0, q_1, \cdots, q_{m}}$,
	\item $\delta$ is the union of the following transitions set:
	\begin{itemize}
		\item $\myset{q_i\xrightarrow[\myvec{0}]{\Sigma}q_i\mid i\in [0, m-1]}$,
		\item $\myset{q_i\xrightarrow[\myvec{0}]{\arrseparator}q_{i+1}\mid i\in[0,m-2]}$,
		\item $\myset{q_{m-1}\xrightarrow[\myvec{0}]{\arrseparator}q_0^x \mid q_0^x\in I_x}$,
		\item $\myset{q\xrightarrow[\myvec{v}]{a}q'\mid q\xrightarrow[\myvec{v}]{a}q'\in \delta_x}$,
		\item $\myset{q_f^x\xrightarrow[\myvec{0}]{\arrseparator}q_m \mid q_f^x\in F_x}$,
		\item $\myset{q_m\xrightarrow[\myvec{0}]{\Sigma\cup\arrseparator} q_m}$,
	\end{itemize}
	\item $I = \myset{q_0}$, $F = \myset{q_m}$, and $\alpha = \alpha_x$.
\end{itemize}

\paragraph{- $A = \mywrite(B, x, m)$:} Suppose that $A$ is constrained by $\aut_A=(R_A, Q_A, \Sigma\cup\{\arrseparator\}, \delta_A, I_A, F_A, \alpha_A)$ and $\aut_A = \aut_{pre}\cdot\aut_{ele}\cdot\aut_{post}$. Then there are $O(|\aut_A|)$ possible concatenation portfolios $(\aut_{pre},\aut_{ele}$, $\aut_{post})$, and for each portfolio satisfying $\aut_{pre}\cap\aut_{m-1}\not=\emptyset$ and $\aut_{ele}\cap\aut_1\not=\emptyset$, we obtain a tuple $((\aut_{pre}\cap\aut_{m-1})\cdot\aut_1\cdot\aut_{post}, \aut_{ele}\cap\aut_1)$ which is one of the pre-image for array write function. The final pre-image is a list of these tuples.

\paragraph{- $x = \myjoin(A,u)$:} Suppose that $x$ is constrained by $\aut_x=(R_x, Q_x, \Sigma, \delta_x, I_x, F_x, \alpha_x)$ and $u=a_1a_2\cdots a_n$, the pre-image for array join function, $\aut_{join}=(R_x, Q_x\cup\myset{q_f}, \Sigma, \delta, I_x, \myset{q_f}, \alpha_x)$, is generated from $\aut_x$ by adding the transition $q\xrightarrow{\arrseparator}q'$ for each $q\xrightarrow{u}q'\in\aut_x$ and $q_f^x\xrightarrow{\arrseparator}q_f$ for each $q_f^x\in F_x$. For $q\xrightarrow{u}q'\in \aut_x$, we mean the state of $\aut_x$ transfer from $q$ to $q'$ after reading the word $u$, that is, there is a path $q\xrightarrow{a_1}q_1\cdots q_{n-1}\xrightarrow{a_n}q'$ in $\aut_x$. 

\paragraph{- $A = \mysplit(x, e)$:} Consider $A$ as string, the semantic of this operation is equal to $A = replaceAll(x, e, \arrseparator)$. We can apply the computation of pre-image of replaceAll function in \cite{ostrich} directly.
}
